\documentclass[../DoAn.tex]{subfiles}
\begin{document}

\begin{center}
    \Large{\textbf{TÓM TẮT NỘI DUNG ĐỒ ÁN}}\\
\end{center}
\vspace{1cm}

Trong bối cảnh số lượng phương tiện giao thông cá nhân ngày càng gia tăng, tình trạng trộm cắp xe máy vẫn đang là một vấn đề đáng quan ngại tại nhiều khu vực, đặc biệt là các đô thị lớn và cao điểm tại thời điểm diễn ra các giải bóng đá lớn như World Cup 2026. Mặc dù hiện nay đã có nhiều giải pháp được áp dụng như khóa cơ học, khóa điện tử, hệ thống báo động chống trộm và thiết bị định vị GPS, các phương pháp này vẫn tồn tại một số hạn chế như khả năng giám sát thời gian thực còn yếu, thiếu tính chủ động trong cảnh báo, hoặc chưa tích hợp đầy đủ giữa chức năng phát hiện, cảnh báo và định vị. Bên cạnh đó, các nghiên cứu gần đây về hệ thống chống trộm dựa trên Internet of Things (IoT) đã mở ra hướng tiếp cận mới khi kết hợp các cảm biến, vi điều khiển và mạng truyền thông để giám sát và phản ứng theo thời gian thực, tuy nhiên vẫn gặp hạn chế về chi phí, độ phức tạp triển khai và độ ổn định của kết nối.

Trong đồ án này, em lựa chọn hướng tiếp cận xây dựng hệ thống chống trộm xe ứng dụng IoT nhằm tận dụng ưu điểm về khả năng kết nối, mở rộng linh hoạt và chi phí thấp của các nền tảng phần cứng phổ biến như ESP32 và các module cảm biến. Hướng tiếp cận này cho phép tích hợp đồng thời các chức năng phát hiện rung động bất thường, xác định vị trí phương tiện thông qua GPS và truyền cảnh báo đến người dùng thông qua mạng GSM 

Giải pháp được xây dựng bao gồm một thiết bị gắn trên xe tích hợp vi điều khiển, cảm biến chuyển động/rung, module định vị và module truyền thông, kết hợp với giao diện giám sát trên thiết bị di động. Khi phát hiện hành vi xâm nhập hoặc di chuyển bất thường, hệ thống sẽ tự động gửi cảnh báo và cung cấp vị trí hiện tại của phương tiện theo thời gian thực.

Đóng góp chính của đồ án là xây dựng được một hệ thống chống trộm xe ứng dụng IoT hoạt động ổn định theo thời gian thực, có khả năng cảnh báo sớm và theo dõi phương tiện từ xa, đồng thời chứng minh tính khả thi của việc áp dụng IoT trong các hệ thống an ninh chi phí thấp. Kết quả đạt được là một hệ thống có thể hoạt động thực tế và có tiềm năng mở rộng trong tương lai.
\begin{flushright}
Sinh viên thực hiện\\
\begin{tabular}{@{}c@{}}
\textit{(Ký và ghi rõ họ tên)}
\end{tabular}
\end{flushright}

\end{document}